\chapter{AI service}

Chương này mô tả AI service ở mức tổ chức: các nhóm dịch vụ, cách quản lý vòng đời mô hình, đường ống giọng nói, các tác tử phụ trợ và danh sách điểm cuối. Bộ não của dịch vụ — pipeline TRACE-CAG — được phân tích riêng và sâu hơn nhiều ở Chương~8.

\section{Vai trò và tổ chức}

AI service là bộ não ngôn ngữ của hệ thống. Nó không sở hữu dữ liệu người học, nhưng
quyết định \textbf{nội dung} mà người học nhận được: câu trả lời của gia sư ảo, chẩn đoán
lỗi, gợi ý luyện tập, bản dịch, phản hồi phát âm và cả nội dung khoá học sinh tự động.

\begin{longtable}{L{4.4cm} L{9.2cm}}
\toprule
\textbf{Nhóm} & \textbf{Thành phần}\\
\midrule
\endhead
Pipeline chính & \code{trace\_cag/} — 19 tệp, 8.078 dòng: đồ thị trạng thái, các đỉnh xử lý, bộ nhớ đệm, truy xuất, sinh phản hồi, đánh giá.\\
Điều phối & \code{orchestrator.py} (khởi tạo và làm ấm mô hình), \code{model\_gateway.py} (vòng đời mô hình), \code{v3\_pipeline.py}.\\
Tri thức & \code{kg\_service\_v3.py} (đồ thị tri thức), \code{jit\_graph\_service.py} (dựng đồ thị con tức thời), \code{subgraph\_hot\_cache.py}, \code{graph\_analytics.py}, \code{kg\_data\_loader.py}.\\
Truy xuất & \code{retrieval\_service\_v3.py}, \code{embedding\_service\_v3.py}, \code{vector\_store\_service.py}, \code{document\_intelligence.py}.\\
Giọng nói & \code{stt/} (21 tệp: nhận dạng luồng, VAD, đóng câu, TTS luồng), \code{tts\_service.py}, \code{hubert\_service.py} (phân tích âm vị).\\
Hội thoại chủ đề & \code{topic\_chat\_service.py}, \code{topic\_catalog\_service.py}, \code{topic\_prompt\_builder.py}, \code{topic\_llm\_gateway.py}, \code{topic\_preloader.py}.\\
Tác tử & \code{content\_agent/} (sinh khoá học), \code{notification\_agent/} (soạn thông báo), \code{ranking\_agent\_insights.py}.\\
Dữ liệu ngoài & \code{content\_etl/} — ETL từ vựng có kiểm soát bản quyền.\\
Trình xử lý mô hình & \code{handlers/}: Gemini, Ollama-Qwen, Qwen, Whisper, Piper, HuBERT, MiniLM.\\
Quan trắc & \code{metrics.py}, \code{telemetry.py}, \code{performance\_monitor.py}, \code{logging\_service.py}.\\
\bottomrule
\caption{Tổ chức AI service}
\end{longtable}

\section{Quản lý mô hình}

\subsection{Model Gateway}

\code{model\_gateway.py} là nơi duy nhất biết mô hình nào đang nằm trong bộ nhớ. Nó cung
cấp: sổ đăng ký mô hình, nạp lười (\en{lazy loading}) có khoá bất đồng bộ, tự động gỡ mô
hình nhàn rỗi, bảng định tuyến theo loại tác vụ, và theo dõi bộ nhớ bằng \code{psutil}.

\begin{longtable}{L{2.6cm} L{2.6cm} L{2.0cm} L{6.0cm}}
\toprule
\textbf{Mô hình} & \textbf{Loại} & \textbf{Bộ nhớ} & \textbf{Dùng cho}\\
\midrule
\endhead
Qwen & hội thoại & \(\approx\)4 GB & Phân tích ngữ pháp, hội thoại\\
Whisper & STT & \(\approx\)1,5 GB & Nhận dạng giọng nói\\
Piper & TTS & \(\approx\)100 MB & Tổng hợp giọng nói\\
HuBERT & phát âm & \(\approx\)2 GB & Phân tích âm vị\\
MiniLM & nhúng & \(\approx\)90 MB & Tìm kiếm ngữ nghĩa\\
Llama-VI & tiếng Việt & \(\approx\)4 GB & Giải thích tiếng Việt\\
\bottomrule
\caption{Sổ đăng ký mô hình}
\end{longtable}

Khi tổng bộ nhớ vượt \code{MAX\_MEMORY\_MB} (mặc định 8 GB), cổng gỡ các mô hình ưu tiên
thấp theo thứ tự lâu không dùng nhất. Vòng đời một mô hình:
\code{Unloaded} $\rightarrow$ \code{Loading} $\rightarrow$ \code{Ready} $\rightleftarrows$
\code{Busy} $\rightarrow$ \code{Unloading}.

\subsection{Chuỗi dự phòng nhà cung cấp}

Hệ thống không phụ thuộc một nhà cung cấp LLM duy nhất. Thứ tự thử: nhà cung cấp đám mây
chính $\rightarrow$ Gemini $\rightarrow$ Ollama chạy cục bộ. Bộ định tuyến thông minh
(\code{smart\_router}) phân loại độ phức tạp câu hỏi để chọn mô hình:

\begin{itemize}[leftmargin=1.4em]
  \item Câu chào hỏi ngắn ($\le$5 từ) $\rightarrow$ mô hình nhỏ chạy cục bộ.
  \item Câu hỏi ngữ pháp, câu dài ($>$50 từ), hoặc chứa thuật ngữ kỹ thuật $\rightarrow$
  mô hình đám mây.
  \item Mặc định $\rightarrow$ mô hình đám mây.
\end{itemize}

\begin{luuy}[Nhà cung cấp gỡ mô hình không báo trước]
Groq đã ngừng \code{llama-3.1-8b-instant} và \code{llama-3.3-70b-versatile} — cả hai đều
từng là mặc định. Mặc định hiện tại là \code{openai/gpt-oss-20b} cho toàn dịch vụ và
\code{openai/gpt-oss-120b} cho tác tử nội dung. Vì mọi đường sinh nội dung đều có
\en{fallback} về khuôn mẫu tất định, một khoá chết trông giống như "sinh thành công"
nhưng nội dung lại là câu mẫu — phải kiểm tra chữ trước khi tin một tác vụ đã chạy đúng.
\end{luuy}

\section{Đường ống giọng nói}

Thư mục \code{stt/} chứa một hệ con hoàn chỉnh gồm 21 tệp, hoạt động theo mô hình
\textbf{song luồng} (\en{dual-stream}): nghe và nói diễn ra đồng thời.

\begin{longtable}{L{3.8cm} L{9.8cm}}
\toprule
\textbf{Thành phần} & \textbf{Vai trò}\\
\midrule
\endhead
\code{audio\_ingest}, \code{ring\_buffer} & Nhận khối âm thanh từ WebSocket vào bộ đệm vòng.\\
\code{vad\_endpointing} & Phát hiện có tiếng nói và xác định điểm kết thúc phát ngôn.\\
\code{engines/}, \code{model\_registry} & Các động cơ nhận dạng (Faster-Whisper và tương đương) cùng sổ đăng ký.\\
\code{sentence\_splitter}, \code{sentence\_finalizer} & Cắt và chốt câu từ dòng bản ghi tạm.\\
\code{transcript\_normalizer}, \code{transcript\_state} & Chuẩn hoá và giữ trạng thái bản ghi.\\
\code{duplex\_turn}, \code{session\_manager} & Quản lý lượt nói hai chiều và phiên; xử lý việc người học ngắt lời AI.\\
\code{streaming\_tts} & Tổng hợp giọng nói theo từng đoạn để giảm độ trễ cảm nhận.\\
\code{verifier\_router} & Định tuyến sang mô hình kiểm chứng khi bản ghi có độ tin cậy thấp.\\
\code{trace\_cag\_adapter} & Chuyển bản ghi đã chốt vào pipeline TRACE-CAG.\\
\code{voice\_ticket} & Vé phiên giọng nói — kiểm soát truy cập WebSocket.\\
\bottomrule
\caption{Thành phần đường ống giọng nói}
\end{longtable}

\subsection{Phân tích phát âm bằng HuBERT}

Âm thanh 16 kHz được đưa qua HuBERT-large để nhận dạng chuỗi âm vị, so sánh với phát âm
mục tiêu và cho điểm từng âm vị. Hệ thống nhận diện riêng các lỗi đặc trưng của người
Việt:

\begin{longtable}{L{3.4cm} L{3.0cm} L{3.0cm} L{4.0cm}}
\toprule
\textbf{Âm vị} & \textbf{Ví dụ} & \textbf{Lỗi thường gặp} & \textbf{Kết quả nghe được}\\
\midrule
\endhead
theta (\en{th} vô thanh) & \en{think} & thay bằng /t/ & "tink"\\
eth (\en{th} hữu thanh) & \en{this} & thay bằng /d/ & "dis"\\
esh (\en{sh}) & \en{shoe} & thay bằng /s/ & "sue"\\
/v/ & \en{very} & thay bằng /b/ & "bery"\\
\bottomrule
\caption{Lỗi phát âm đặc trưng của người học Việt Nam}
\end{longtable}

\section{Các tác tử AI khác}

\begin{longtable}{L{3.4cm} L{10.2cm}}
\toprule
\textbf{Tác tử} & \textbf{Chức năng}\\
\midrule
\endhead
\code{content\_agent} & Sinh khoá học hoàn chỉnh: \code{planner} lập chương trình theo bậc CEFR, \code{generator} sinh bài học và bài tập, \code{policies} ràng buộc hình dạng, \code{adapters} nối các nguồn từ vựng, \code{store} lưu kết quả. Backend \code{content\_agent\_apply} mới là nơi ghi vào cơ sở dữ liệu.\\
\code{notification\_agent} & Đồ thị trạng thái riêng (\code{graph}, \code{nodes}, \code{state}) soạn nội dung thông báo cá nhân hoá.\\
\code{ranking\_agent\_insights} & Sinh nhận định về vị trí xếp hạng và gợi ý cải thiện.\\
\code{error\_diagnosis} & Điểm cuối nội bộ chẩn đoán lỗi độc lập, dùng cho sổ tay lỗi sai.\\
\code{translate} & Dịch có ngữ cảnh học tập, khác dịch máy thuần tuý ở chỗ giữ mức độ khó phù hợp bậc CEFR.\\
\code{topic\_chat} & Hội thoại theo chủ đề với lời nhắc và danh mục riêng, nhẹ hơn TRACE-CAG đầy đủ.\\
\bottomrule
\caption{Các tác tử AI phụ trợ}
\end{longtable}

\section{Điểm cuối AI service}

\begin{longtable}{L{4.6cm} L{9.0cm}}
\toprule
\textbf{Tiền tố} & \textbf{Chức năng}\\
\midrule
\endhead
\code{/api/v1/chat} & Phiên hội thoại cơ bản.\\
\code{/api/v1/lexi*} & Trợ lý Lexi chạy TRACE-CAG đầy đủ, hỗ trợ dòng SSE.\\
\code{/api/v1/stt} & Nhận dạng giọng nói và chấm điểm phát âm.\\
\code{/api/v1/tts} & Tổng hợp giọng nói.\\
\code{/api/v1/voice} & Phiên giọng nói thời gian thực (WebSocket).\\
\code{/api/v1/topics} & Hội thoại theo chủ đề.\\
\code{/api/v1/ai} & Phân tích AI, tình trạng mô hình, dịch.\\
\code{/api/v1/admin} & Quản trị mô hình.\\
\code{/warmup} & Kích hoạt nạp trước mô hình lúc triển khai.\\
\code{/visualizer} & Trực quan hoá các đỉnh của pipeline.\\
\bottomrule
\caption{Điểm cuối AI service}
\end{longtable}

\begin{luuy}[Làm ấm khi khởi động, không phải khi có yêu cầu đầu tiên]
Trước đây bộ điều phối và mô hình Qwen chỉ được nạp khi có yêu cầu đầu tiên, khiến người
dùng đầu tiên sau mỗi lần triển khai phải chờ rất lâu. Việc nạp nay nằm trong hàm
\code{lifespan()} của \code{main.py}, cộng với việc song song hoá tra cứu MongoDB/Redis
trong \code{chat.py}. Riêng \code{RetrievalServiceV3} mất khoảng 169,5 giây để dựng đồ thị
NetworkX và tính độ trung tâm trên đồ thị sản xuất, nên được bọc trong
\code{asyncio.to\_thread} kèm khoá để không chặn vòng lặp sự kiện và không bị dựng hai lần.
\end{luuy}
